\chapter{Problem Understanding and Positioning}

\section{What the CPI is, precisely}

The Consumer Price Index measures the change over time in the general level of retail
prices of a fixed basket of goods and services that households buy for consumption. It is
India's headline retail inflation measure and the operational target of the RBI's flexible
inflation targeting framework. Understanding four technical terms is a prerequisite for
everything in Part II.

\begin{center}
\small
\begin{tabularx}{\textwidth}{@{}L{4.2cm} X@{}}
\toprule
\bhead{Term} & \bhead{Meaning, and its value in CPI 2024} \\
\midrule
\textbf{Index reference period} & The period for which the index equals 100. For the
current series: calendar year \textbf{2024 $=$ 100}. \\
\addlinespace[2pt]
\textbf{Weight reference period} & The period from which expenditure weights are derived.
For CPI 2024: \textbf{HCES 2023--24}, using the Mixed Modified Reference Period. \\
\addlinespace[2pt]
\textbf{Price reference period} & The period during which base prices are collected.
For CPI 2024: \textbf{calendar year 2024}. \\
\addlinespace[2pt]
\textbf{Elementary aggregate} & The lowest level at which an index is computed, and the
only level at which no expenditure weights are available. This is where the choice of
Jevons vs.\ Dutot vs.\ Carli lives, and it is where APIx does most of its work. \\
\bottomrule
\end{tabularx}
\end{center}

\subsection{The CPI 2024 series: the facts you must have straight}

\begin{center}
\small
\begin{tabularx}{\textwidth}{@{}L{6.2cm} X@{}}
\toprule
\bhead{Attribute} & \bhead{CPI 2024 (base 2024$=$100)} \\
\midrule
First release & 12 February 2026, for January 2026 (2.75\% y-o-y, provisional) \\
Predecessor & CPI 2012 series, last released for December 2025 \\
Weights source & HCES 2023--24 (MMRP) \\
Classification & \textbf{COICOP 2018} (UN Statistics Division): 12 Divisions, 43 Groups, 92 Classes, 162 Sub-classes \\
Weighted items & 358 at All-India level (308 goods, 50 services), up from 299 \\
Sample & 1{,}465 rural markets, 1{,}395 urban markets across 434 towns, \textbf{plus 12 online markets} \\
Collection & Tablet-based CAPI by the Field Operations Division of NSO; online prices weekly \\
Elementary formula & \textbf{Jevons} \\
Higher-level formula & \textbf{Young / Modified Laspeyres} \\
Airfare collection & \textbf{``through well-known online platforms''} \\
Rail fare & Administrative data from the Ministry of Railways \\
Fuel & Retail selling prices from PPAC, Ministry of Petroleum \& Natural Gas \\
Dissemination & e-Sankhyiki portal; monthly indices for Divisions, Groups, Classes, Sub-classes and items, All-India and State/UT \\
Back series & Jan 2013 -- Dec 2024, computed with a linking factor over overlap year 2025 \\
Linking factors (General) & Rural 0.5222, Urban 0.5320, Combined 0.5267 \\
\bottomrule
\end{tabularx}
\end{center}
\srcnote{MoSPI, CPI 2024 FAQ (Annexure V, 12 Feb 2026); MoSPI CPI press releases Jan--Jun 2026; MoSPI Expert Group Report on Comprehensive Updation of CPI.}

\subsection{Where air fare sits in the tree}

Under COICOP 2018 the relevant path is:

\begin{center}
\small
\begin{tabular}{@{}l l l@{}}
\toprule
\bhead{Level} & \bhead{Code} & \bhead{Label} \\
\midrule
Division & 07 & Transport \\
Group & 07.3 & Transport services \\
Class & 07.3.3 & Passenger transport by air \\
Sub-class & 07.3.3.x & Domestic / international air passenger transport \\
\bottomrule
\end{tabular}
\end{center}

Division-level weights in CPI 2024 are published:

\begin{center}
\small
\begin{tabular}{@{}l r r r r r r@{}}
\toprule
& \multicolumn{2}{c}{\bhead{Rural}} & \multicolumn{2}{c}{\bhead{Urban}} & \multicolumn{2}{c}{\bhead{Combined}} \\
\cmidrule(lr){2-3}\cmidrule(lr){4-5}\cmidrule(lr){6-7}
\bhead{Division} & 2012 & 2024 & 2012 & 2024 & 2012 & 2024 \\
\midrule
Food and beverages & 50.922 & 41.983 & 32.811 & 30.251 & 42.617 & 36.753 \\
Housing, water, electricity, gas & 7.983 & 11.764 & 27.294 & 25.000 & 16.888 & 17.665 \\
\rowcolor{lightband} \textbf{Transport (07)} & \textbf{5.645} & \textbf{8.644} & \textbf{7.129} & \textbf{8.985} & \textbf{6.394} & \textbf{8.796} \\
Information and communication & 2.818 & 3.647 & 3.906 & 3.563 & 3.323 & 3.609 \\
Health & 6.839 & 6.764 & 4.820 & 5.275 & 5.900 & 6.100 \\
\bottomrule
\end{tabular}
\end{center}
\srcnote{MoSPI CPI 2024 FAQ, Q.39. Abridged --- the full 12-division table is in the source.}

\begin{keybox}
\bhead{[VERIFY] before the pitch.} The \emph{item-level} weight for passenger transport by
air is published in the item-wise weight annexure at
\url{https://cpi.mospi.gov.in} (Announcements tab) and on \url{https://mospi.gov.in}.
Pull the exact number the week of the hackathon and quote it. Two things to note when you
do: (i) the Transport division rose from 6.394 to 8.796 combined --- a 38\% relative
increase in weight, which is a strong ``this matters more now'' talking point; (ii) air
fare is only a slice of 07.3, so do not overstate the headline impact. Honest framing
beats a big number.
\end{keybox}

\section{The institutional map: who consumes what}

A pitch that names the right desks is worth more than one that says ``the government''.

\begin{center}
\small
\begin{tabularx}{\textwidth}{@{}L{3.6cm} L{3.6cm} X@{}}
\toprule
\bhead{Body} & \bhead{Unit} & \bhead{Relationship to APIx} \\
\midrule
MoSPI / NSO & \textbf{Price Statistics Division (PSD)} & Owns CPI methodology, weights, item basket, COICOP mapping. \textbf{This is who the PS means by ``PSD''} in ``index-construction module based on PSD given routes and weights.'' PSD is the consumer of the index engine and the supplier of the weight vector. \\
\addlinespace[2pt]
MoSPI / NSO & Field Operations Division (FOD) & Runs the actual price collection, CAPI tablets, 434 towns. APIx does not replace FOD; it supplements the online-market stream. \\
\addlinespace[2pt]
MoSPI & \textbf{Data Informatics \& Innovation Division (DIID)} & \textbf{The PS owner.} Mandate is alternative data sources and modernisation. Speak to DIID's incentives: reusable pipeline, not a one-off scraper. \\
\addlinespace[2pt]
RBI & Monetary Policy Department & Consumes CPI for the flexible inflation targeting framework. Values \emph{timeliness} above all --- a credible T$+$1 nowcast of a division is directly useful. \\
\addlinespace[2pt]
MoCA / DGCA & \textbf{Tariff Monitoring Unit (TMU)} & Established 2010. Monitors fares on \textbf{78 selected routes}, \emph{monthly}, \emph{manually}, \emph{by looking at airline websites}, covering roughly \textbf{27\% of domestic traffic}. This is the unit whose work APIx most directly automates. \\
\addlinespace[2pt]
AERA & Tariff orders & Sets User Development Fee per airport. Source of truth for the UDF component of fare decomposition. \\
\addlinespace[2pt]
BCAS / MoCA & ASF notifications & Aviation Security Fee. Second statutory levy component. \\
\bottomrule
\end{tabularx}
\end{center}
\srcnote{Rajya Sabha reply by MoS Civil Aviation, 8 Dec 2025 (TMU, 78 routes, $\approx$27\% of traffic); MoSPI Expert Group Report (PSD role in COICOP mapping).}

\subsection{The DGCA TMU is the strongest single hook in this PS}

Three facts, in sequence, are the most persuasive thirty seconds you have:

\begin{enumerate}
\item The DGCA's Tariff Monitoring Unit monitors \textbf{78 routes, manually, once a
month, by opening airline websites} --- covering about 27\% of domestic traffic.
\item In its 375th Report (\emph{Demands for Grants 2025--26, Ministry of Civil
Aviation}), the Parliamentary Standing Committee on Transport, Tourism and Culture found
that the DGCA \emph{lacks capacity to proactively monitor airfare trends}, and recommended
an AI-driven predictive monitoring system, a public \textbf{``Pricing Transparency
Index''}, and an \textbf{``Airfare Vigil''} reporting app.
\item In July 2026 the Ministry told Parliament that average airfare on \textbf{72 domestic
sectors rose 20.5\%} comparing June 2026 with March 2025, attributing it to ATF --- which
is 35--40\% of airline operating cost and which spiked over 100\% on the West Asia
conflict.
\end{enumerate}

\noindent APIx is, in effect, a working prototype of two out of three of those
Parliamentary recommendations, built for the statistics ministry rather than the aviation
ministry. That is a genuinely strong positioning line and almost no competing team will
find it.

\section{Why airfare is the hardest item in the CPI basket}

\begin{center}
\small
\begin{tabularx}{\textwidth}{@{}L{4.3cm} X@{}}
\toprule
\bhead{Pathology} & \bhead{Why it breaks naive price collection} \\
\midrule
\textbf{Dynamic pricing} & The same seat on the same flight has a different price every
hour. There is no ``the price'' to observe --- only a price \emph{conditional on when you
look and how far ahead you are looking}. \\
\addlinespace[2pt]
\textbf{Perishable inventory} & A seat unsold at departure is worthless. Revenue management
therefore raises price as inventory depletes and as departure approaches --- and sometimes
\emph{cuts} it to dump distressed inventory. The lead-time--price relation is not monotone. \\
\addlinespace[2pt]
\textbf{Bucketing} & Airlines sell from discrete fare buckets (RBDs). Observed prices are a
step function, not a continuum, so the empirical distribution is multi-modal. Mean-based
statistics are badly behaved; log-space and rank-based statistics are not. \\
\addlinespace[2pt]
\textbf{Availability censoring} & When the cheap buckets sell out you observe only
expensive quotes. If you drop the missing cheap ones as ``no data'', the index is
\emph{biased downward} in exactly the periods you care about (festivals, disruptions). This
is the classic selection-on-the-dependent-variable trap. \\
\addlinespace[2pt]
\textbf{Bundle drift} & ``Base fare'' means different things across carriers: cabin bag
only, checked bag included, seat selection, meal, change fee. Comparing an IndiGo
\emph{Saver} with an Air India \emph{Comfort} is a quality comparison, not a price
comparison. \\
\addlinespace[2pt]
\textbf{Personalisation} & Prices may vary by cookie state, device, geography and account.
An unauthenticated, cookie-clean, India-egress collection is a \emph{methodological
requirement}, not a convenience. \\
\addlinespace[2pt]
\textbf{Route entry and exit} & Carriers add and drop sectors constantly. A fixed matched
sample decays; an unmatched sample chain-drifts. This is precisely the problem multilateral
index methods were invented to solve. \\
\addlinespace[2pt]
\textbf{Acquisition vs.\ use} & Do you record the price when the ticket is bought or when
the flight is taken? The two give different indices. ONS records the change \textbf{in the
month the flight departs, not when the ticket is bought.} Most teams will not even know
this is a question. \\
\bottomrule
\end{tabularx}
\end{center}

\section{The Indian market, in numbers you can quote}

\begin{center}
\small
\begin{tabularx}{\textwidth}{@{}L{7.2cm} X@{}}
\toprule
\bhead{Fact} & \bhead{Source / date} \\
\midrule
Domestic airlines carried \textbf{864.04 lakh} passengers in Jan--Jun 2026, $+$1.44\% y-o-y & DGCA traffic report, June 2026 \\
July 2026: $\approx$1.20 crore domestic passengers, $-$4.8\% y-o-y; Jan--Jul 2026 total 984.03 lakh & DGCA, Aug 2026 \\
IndiGo domestic market share \textbf{67.4\%} (July 2026); record 66.3\% in June 2026 & DGCA monthly \\
Air India Group $\approx$24\%; Akasa 5.5\%; SpiceJet 1.6\% (July 2026) & DGCA monthly \\
Passenger load factors June 2026: Akasa 92.2\%, SpiceJet 87.8\%, AI Group 85.5\%, IndiGo 85.1\% & DGCA monthly \\
Cancellation rate 0.63\% (June 2026); technical 39.6\%, operational 36.3\%, weather 20.3\% & DGCA monthly \\
ATF is \textbf{35--40\% of airline operating cost}; ATF spiked $>$100\% on the West Asia conflict & MoCA reply to Parliament, July 2026 \\
Average fare on \textbf{72 domestic sectors up 20.5\%}, June 2026 vs.\ March 2025 & MoCA reply to Parliament, July 2026 \\
DGCA TMU monitors \textbf{78 routes} manually, monthly, $\approx$27\% of domestic traffic & Rajya Sabha, 8 Dec 2025 \\
India: third-largest civil aviation market, $\approx$4.2\% of global air traffic (2024) & IATA / industry \\
\bottomrule
\end{tabularx}
\end{center}

\begin{warnbox}
\bhead{Market-structure consequence for the index design.} With one carrier at $\approx$67\%
share, a carrier-equal-weighted index is \textbf{wrong}. It would let SpiceJet's 1.6\%-share
pricing move the national index as much as IndiGo's 67\%. APIx therefore weights carriers
\emph{within} each route by that carrier's share of passengers on that route, not by
national share and not equally. This is a detail that separates a statistics project from
a dashboard project --- see \S\ref{sec:carrierweights}.
\end{warnbox}

\section{Gap analysis: what genuinely does not exist today}

\begin{center}
\small
\begin{tabularx}{\textwidth}{@{}L{4.2cm} C{2.3cm} C{2.3cm} X@{}}
\toprule
\bhead{Capability} & \bhead{CPI 2024} & \bhead{DGCA TMU} & \bhead{APIx} \\
\midrule
Automated collection & Partially & No (manual) & Yes \\
Daily frequency & No (weekly collect) & No (monthly) & Yes \\
Route-level index & No & Fare levels only & Yes \\
Advance-purchase dimension & No & No & Yes, 5 windows \\
Availability adjustment & No & No & Yes \\
Fare-component decomposition & No & No & Yes \\
Published elasticity surface & No & No & Yes \\
Machine-readable public API & Portal only & No & Yes (SDMX-JSON) \\
Provenance audit trail & Internal & Internal & Public endpoint \\
Nowcast of Division 07 & No & No & Yes \\
\bottomrule
\end{tabularx}
\end{center}

\noindent The honest conclusion: \textbf{APIx is not filling a hole in CPI; it is building
the instrument layer that neither the statistics ministry nor the aviation regulator
currently has, and handing the same instrument to both.} That is the sentence to put on the
final slide.
